\chapter{Arc-Length Comparison: Parabola vs.\ Helix}\label{ch:examples}

\section{Principle of the Comparison}

The quadratic curve $k(n) = \frac{1}{6}(2n^2 + 3n + 1)$ was geometrically
wound onto a cone, where it forms the helix $\gamma(\theta)$.
The central question is: how closely do the arc lengths agree?


\textbf{Arc length of the parabola} between $n = n_1$ and $n = n_2$:
\[
L_{\text{Parabola}} = \int_{n_1}^{n_2} \sqrt{1 + \left(\frac{dk}{dn}\right)^2}\, dn
= \int_{n_1}^{n_2} \frac{1}{6}\sqrt{36 + (4n+3)^2}\, dn.
\]

\textbf{Arc length of the helix} with the direct $\theta$-parametrisation:
\[
\gamma(\theta) = \begin{pmatrix}
\frac{6}{\pi^2}\theta\cos\theta\\[4pt]
\frac{6}{\pi^2}\theta\sin\theta\\[4pt]
-\frac{3}{4} + \frac{3}{\pi}\theta
\end{pmatrix},
\qquad
\theta_n = \frac{\pi}{3}n + \frac{\pi}{4}.
\]

Differentiation yields:
\[
\|\gamma'(\theta)\| = \sqrt{\left(\frac{6}{\pi^2}\right)^2(1 + \theta^2)
+ \left(\frac{3}{\pi}\right)^2}.
\]

The arc length is therefore:
\[
L_{\text{Helix}} = \int_{\theta_{n_1}}^{\theta_{n_2}} \|\gamma'(\theta)\|\, d\theta.
\]

\section{Numerical Examples}

All arc lengths were computed numerically to high precision.
As a detailed example, the first step
($n = 1 \to n = 5$, major step / gap zone) is worked out step by step;
the remaining results are summarised in
Table~\ref{tab:arc-length-comparison}.

% ═══════════════════════════════════════════════════════════════════════
\subsection{Detailed Example: \texorpdfstring{$n = 1 \rightarrow n = 5$}{n = 1 → n = 5}
 (Major Step)}
% ═══════════════════════════════════════════════════════════════════════

\textbf{Step~1 --- Determine parameters:}
\[
n_1 = 1,\quad n_2 = 5,\quad \Delta n = 4,\quad
\theta_1 = \frac{\pi}{3}\cdot 1 + \frac{\pi}{4} = 1{,}8326,\quad
\theta_2 = \frac{\pi}{3}\cdot 5 + \frac{\pi}{4} = 6{,}0214.
\]

\textbf{Step~2 --- Arc length of the parabola:}
\[
L_{\text{Parabola}}
= \int_{1}^{5} \frac{1}{6}\sqrt{36 + (4n+3)^2}\, dn
= \boxed{10{,}839409}.
\]

\textbf{Step~3 --- Arc length of the helix:}
\[
L_{\text{Helix}}
= \int_{1{,}8326}^{6{,}0214}
  \sqrt{\left(\frac{6}{\pi^2}\right)^2(1 + \theta^2)
  + \left(\frac{3}{\pi}\right)^2}\, d\theta
= \boxed{11{,}155199}.
\]

\textbf{Step~4 --- Relative deviation:}
\[
\varepsilon = \frac{L_{\text{Helix}} - L_{\text{Parabola}}}{L_{\text{Parabola}}}
= \frac{0{,}316}{10{,}839} = \boxed{+2{,}91\,\%}.
\]

\begin{remark}\label{rem:helix-longer}
The helix is \emph{longer} than the parabola, since the spiral motion in
the $(x,y)$-plane (especially at small $\theta$) generates additional
path length.
\end{remark}

% ═══════════════════════════════════════════════════════════════════════
\subsection{Comparison Table}
% ═══════════════════════════════════════════════════════════════════════

\begin{table}[H]
\centering
\caption{Arc-length comparison: parabola vs.\ helix
  (true $\theta$-parametrisation).}
\label{tab:arc-length-comparison}
\begin{tabular}{clccccr}
\toprule
No. & $n_1 \to n_2$ & Type
  & $L_{\text{Parabola}}$ & $L_{\text{Helix}}$
  & Diff. & Dev. \\
\midrule
1 & $1 \to 5$   & Major
  & 10{,}839409 & 11{,}155199 & $+0{,}316$ & $+2{,}91\,\%$ \\
2 & $5 \to 7$   & Minor
  & 9{,}221075  & 9{,}309137  & $+0{,}088$ & $+0{,}96\,\%$ \\
3 & $11 \to 13$ & Minor
  & 17{,}117480 & 17{,}164862 & $+0{,}047$ & $+0{,}28\,\%$ \\
4 & $17 \to 19$ & Minor
  & 25{,}079948 & 25{,}112276 & $+0{,}032$ & $+0{,}13\,\%$ \\
5 & $23 \to 25$ & Minor
  & 33{,}060583 & 33{,}085105 & $+0{,}025$ & $+0{,}07\,\%$ \\
6 & $35 \to 37$ & Minor
  & 49{,}040809 & 49{,}057339 & $+0{,}017$ & $+0{,}03\,\%$ \\
\bottomrule
\end{tabular}
\end{table}

% ═══════════════════════════════════════════════════════════════════════
\section{Unified Representation via the Universal Quantity \texorpdfstring{$Q$}{Q}}
% ═══════════════════════════════════════════════════════════════════════

The convergence of the arc lengths can be explained analytically by
expressing both integrals in terms of the universal quantity
$Q(k) = \sqrt{48k+1} = 4n+3$. With $dQ = 4\,dn$ and
$\theta = \pi Q/12$, hence $d\theta = (\pi/12)\,dQ$, one obtains:

\begin{theorem}[$Q$-Parametrisation of the Arc Lengths]\label{thm:q-arc-length}
The arc-length elements of the parabola $k(n)$ and of the conical helix
can be expressed uniformly as
\begin{equation}\label{eq:ds-uniform}
  ds = \frac{1}{24}\sqrt{Q^2 + C}\; dQ
\end{equation}
where
\[
  C_{\mathrm{Parabola}} = 36,
  \qquad
  C_{\mathrm{Helix}} = 36 + \frac{144}{\pi^2} \approx 50{,}59.
\]
\end{theorem}

\begin{proof}
\textbf{Parabola.}\; The arc-length element of the curve $(n, k(n))$ is
$ds = \sqrt{1 + k'(n)^2}\,dn$ with $k'(n) = (4n+3)/6$.
Substituting $Q = 4n+3$, $dn = dQ/4$:
\[
  ds = \sqrt{1 + \frac{Q^2}{36}}\;\frac{dQ}{4}
     = \frac{1}{4}\cdot\frac{\sqrt{36 + Q^2}}{6}\;dQ
     = \frac{1}{24}\sqrt{Q^2 + 36}\;dQ.
\]

\textbf{Helix.}\; With $a = 6/\pi^2$, $c_1 = 3/\pi$, and
$\theta = \pi Q/12$ we have $d\theta = (\pi/12)\,dQ$. The arc-length
element is $ds = \sqrt{a^2(1+\theta^2) + c_1^2}\,d\theta$. Substituting
$\theta = \pi Q/12$:
\begin{align*}
  a^2(1+\theta^2) + c_1^2
  &= \frac{36}{\pi^4}\Bigl(1 + \frac{\pi^2 Q^2}{144}\Bigr) + \frac{9}{\pi^2}\\
  &= \frac{36}{\pi^4} + \frac{Q^2}{4\pi^2} + \frac{9}{\pi^2}
   = \frac{1}{4\pi^2}\Bigl(Q^2 + 36 + \frac{144}{\pi^2}\Bigr).
\end{align*}
With $d\theta = (\pi/12)\,dQ$ it follows that:
\[
  ds = \frac{1}{2\pi}\sqrt{Q^2 + 36 + \frac{144}{\pi^2}}
       \cdot \frac{\pi}{12}\,dQ
     = \frac{1}{24}\sqrt{Q^2 + 36 + \frac{144}{\pi^2}}\;dQ.
     \qedhere
\]
\end{proof}

\begin{remark}[Geometric Meaning]\label{rem:geometric-meaning-Q}
The difference $C_{\mathrm{Helix}} - C_{\mathrm{Parabola}} = 144/\pi^2
\approx 14{,}59$ is the exact contribution of the spiral motion in the
$(x,y)$-plane. For $Q \to \infty$ the term $Q^2$ dominates in both
cases, and the ratio of the arc-length elements converges:
\[
  \frac{ds_{\mathrm{Helix}}}{ds_{\mathrm{Parabola}}}
  = \sqrt{\frac{Q^2 + 36 + 144/\pi^2}{Q^2 + 36}}
  \;\xrightarrow{Q\to\infty}\; 1.
\]
This explains the convergence $\varepsilon \to 0$ observed in
Table~\ref{tab:arc-length-comparison} analytically, not merely
numerically.
\end{remark}

% ═══════════════════════════════════════════════════════════════════════
\section{Interpretation of the Results}
% ═══════════════════════════════════════════════════════════════════════

\begin{proposition}[Asymptotic Arc-Length Convergence]
\label{prop:arc-length-convergence}
The helix arc length is consistently \emph{longer} than the arc length
of the corresponding parabola segments, with the relative difference
converging monotonically to zero:
\[
\frac{L_{\mathrm{Helix}}(n, n+2) - L_{\mathrm{Parabola}}(n, n+2)}
{L_{\mathrm{Parabola}}(n, n+2)}
\xrightarrow{n \to \infty} 0.
\]
This follows directly from the uniform $Q$-representation
(Theorem~\ref{thm:q-arc-length}), since
$C_{\mathrm{Helix}}/Q^2 \to 0$ and
$C_{\mathrm{Parabola}}/Q^2 \to 0$ as $Q \to \infty$.
\end{proposition}

\begin{corollary}[Practical Accuracy]\label{cor:practical-accuracy}
For all primes $p > 11$ (corresponding to $n \ge 11$) the deviation is
below $0{,}3\,\%$. For applications in the visualisation of prime
distributions this accuracy is more than sufficient.
\end{corollary}
